% ============================================================
% LITERATURE REVIEW UPDATE — Critique DPO
% New papers (2024–2026) to incorporate into Section 3
% ============================================================
% This file can be \input{} into main.tex, or merged manually.
% New \addbibresource entries are collected at the end of this file.
% ============================================================

\section*{Literature Review Update (2024--2026)}
\addcontentsline{toc}{section}{Literature Review Update}

The literature on critique-augmented preference optimisation and
generative verifiers has grown rapidly since this document was first
drafted.  The additions below are organised into five new clusters
that extend the existing review in Section~\ref{sec:lit}.

% ============================================================
\subsection{Step-Level Preference Optimisation}
\label{sec:step-dpo}
% ============================================================

Standard DPO assigns a single preference signal to a complete
response.  A growing line of work argues that CoT reasoning requires
\emph{per-step} supervision because errors early in a chain
propagate forward and cannot be corrected by an outcome-level signal.

% \subsubsection{From r to Q*: DPO as Token-Level Q-Learning}

% \autocite{rafailov2024rqlanguagemodel} prove that the DPO optimum
% satisfies the Bellman equation at the \emph{token} level, so a
% DPO-trained policy is implicitly a Q-function over the token-level
% MDP.  The key identity is
% \begin{equation}
%   \pi^*(y_t \mid x, y_{<t})
%   \;\propto\;
%   \pi_{\mathrm{ref}}(y_t \mid x, y_{<t})
%   \exp\!\bigl(\beta^{-1} A^*(s_t, a_t)\bigr),
% \end{equation}
% where $A^*$ is the optimal advantage function.  This result provides
% the theoretical foundation for step-level DPO variants: because the
% DPO loss already assigns implicit step rewards, supervising
% \emph{intermediate} steps directly (rather than whole responses)
% should improve credit assignment.

% \subsubsection{Chain of Preference Optimisation (CPO)}

% \autocite{zhang2024chainpreferenceoptimization} observe that
% Tree-of-Thought (ToT) search trees naturally expose step-level
% preference pairs: at each branching node one sub-tree eventually
% reaches a correct answer while others do not.  CPO extracts these
% pairs and applies DPO \emph{step by step}, enabling an LLM to
% reproduce ToT quality at inference without search.  Experiments on
% question answering, fact verification, and arithmetic yield up to a
% 4.3\% average improvement over the SFT baseline.

\subsubsection{Step-DPO}

\autocite{lai2024stepdpostepwisepreferenceoptimization}
construct 10K step-wise preference pairs using self-generated data
and apply DPO at the step level, arguing that whole-response
comparison provides insufficient localisation of errors in
long-chain mathematical reasoning.  Applied to Qwen2-72B-Instruct,
Step-DPO achieves 70.8\% on MATH, surpassing GPT-4 at the time of
publication.

% \subsubsection{Step-Controlled DPO (SCDPO)}

% \autocite{lu2024stepcontrolleddpoleveragingstepwise} create negative
% samples for DPO by generating reasoning traces that begin making
% errors at a specified step, explicitly teaching the model to identify
% and avoid reasoning errors at known positions.  Applied to
% InternLM2-20B, SCDPO achieves 88.5\% on GSM8K and 58.1\% on MATH.

% \subsubsection{Self-Explore}

% \autocite{hwang2024selfexporeenhancingmathematical} train LLMs to
% identify the \emph{first pit} (first erroneous step) in their own
% reasoning traces and then apply DPO on step-level preference pairs
% derived from this analysis.  The first-pit identification acts as a
% form of automated self-critique, improving GSM8K and MATH performance
% by 11.6\% and 2.9\% respectively across three base models.  This is
% among the closest existing approaches to the Critique DPO idea: the
% model generates a localised critique before constructing preference
% data.

\subsubsection{Thinking Preference Optimisation (ThinkPO)}

\autocite{yang2025thinkingpreferenceoptimization} apply DPO to
reasoning depth itself, using short CoT responses as rejected
completions and long CoT responses as chosen completions.  Without
any new training data, ThinkPO improves DeepSeek-R1-Distill-Qwen-7B
from 87.4\% to 91.2\% on MATH500 and increases average output length
by 25.9\%.

\subsubsection{Uncovering the Role of CoT in DPO}

\autocite{gao2025uncoveringimpactchainofthought} provide empirical
evidence that DPO's success on reasoning tasks is largely
\emph{attributable} to the presence of CoT data.  Across ten LLMs
and the text-to-SQL domain (which lacks CoT), DPO degrades
performance in six cases.  The analysis shows that CoT in DPO
training mitigates reward hacking and improves discriminative
calibration—directly motivating the Critique DPO approach of
incorporating verifier CoT explicitly into the loss.

% \subsubsection{Advancing LLM Reasoning with Preference Trees (UltraInteract)}

% \autocite{yuan2024eurus} introduce UltraInteract, a large-scale
% dataset structured as preference trees where each instruction has
% multiple correct and incorrect reasoning paths with per-step
% comparison data.  A notable finding is that standard DPO
% \emph{hurts} performance on reasoning tasks, while KTO and NCA
% improve it—further motivating principled modifications to the DPO
% objective such as Critique DPO.

% ============================================================
\subsection{Process Reward Models: From Scalar to Generative}
\label{sec:prm}
% ============================================================

Process Reward Models (PRMs) assign a reward to \emph{each} step of
a reasoning chain rather than to the final answer alone.  The field
has evolved rapidly from human-annotated scalar PRMs toward automated
and generative variants.

\subsubsection{Let's Verify Step by Step (PRM800K)}

\autocite{lightman2023letsverifystepstep} demonstrate that
process supervision substantially outperforms outcome supervision on
MATH, achieving 78\% accuracy with a process-supervised reward model
compared to 72.4\% for outcome supervision.  The paper releases
PRM800K, a dataset of 800K step-level human annotations, and remains
the canonical reference for step-level feedback in reasoning.

% \subsubsection{Automated PRMs: Math-Shepherd and OmegaPRM}

% \autocite{wang2024mathshepherd} introduce an automated process for
% assigning step-level reward labels using Monte-Carlo completion
% estimates, eliminating human annotation.  The step-level PRM is then
% used for both best-of-$n$ reranking and PPO training.

% \autocite{luo2024improvemathematicalreasoninglanguage} (OmegaPRM)
% refine this with a divide-and-conquer MCTS strategy that locates the
% first error in a CoT via binary search, collecting over 1.5M process
% annotations at significantly lower compute cost.  The Google DeepMind
% model improves from 51\% to 69.4\% on MATH500.

% \subsubsection{Process Advantage Verifiers (PAVs)}

% \autocite{setlur2024rewardingprogressscalingautomated} introduce
% Process Advantage Verifiers that measure the \emph{progress} made by
% a step as the change in expected outcome reward before and after the
% step.  PAVs outperform ORMs by over 8\% during test-time search and
% enable 5--6$\times$ sample efficiency gains in online RL, providing a
% principled notion of step value directly relevant to Critique DPO's
% step-level credit assignment.

% \subsubsection{Free Process Rewards and PRIME}

% \autocite{yuan2024freeprocessrewardsprocess} show that implicit PRMs
% can be derived from outcome reward models \emph{for free}, using
% log-likelihood ratios without any step-level labels, outperforming
% Math-Shepherd while using $<$1/38 of its training data.

% \autocite{cui2025primereinforcementprocessimplicit} (PRIME) apply
% these implicit process rewards in online RL (GRPO, PPO, RLOO),
% achieving a 15.1\% average improvement starting from Qwen2.5-Math-7B
% using only 10\% of the data needed by supervised approaches.

\subsubsection{ThinkPRM: Process Reward Models That Think}

\autocite{khalifa2025thinkprmprocessrewardmodels} is perhaps the
closest existing work to the Critique DPO generative-verifier idea at
the process level.  ThinkPRM generates a verification
chain-of-thought \emph{before} scoring each reasoning step rather
than predicting a scalar.  Using only 1\% of the PRM800K labels,
ThinkPRM surpasses discriminative PRMs trained on the full dataset by
8\% on GPQA-Diamond, demonstrating that generative CoT verifiers are
substantially more data-efficient than scalar counterparts.

\subsubsection{GenPRM}

\autocite{zhao2025genprm} extend the generative PRM idea by combining
natural-language reasoning with executable code verification before
assigning a step score.  A 7B GenPRM exceeds Qwen2.5-Math-PRM-72B
using only 23K training examples, further establishing that
generative verifier reasoning provides a strong inductive bias for
process supervision.

\subsubsection{Lessons in Developing PRMs}

\autocite{zhang2025lessonsdevelopingprocessreward} provide an
important empirical caveat: Monte-Carlo estimation methods (e.g.\
Math-Shepherd) yield inferior PRMs compared to LLM-as-judge and
human annotation.  PRMs trained with MC estimation tend to focus
assessment on final-answer steps rather than intermediate steps,
suggesting that critique-style verification (as in GenRM-CoT and
Critique DPO) may better capture the true step-level quality signal.

% \subsubsection{MCTS-Based Self-Training with Process Rewards}

% \autocite{zhang2024restmcts} (ReST-MCTS$^*$) integrate process
% reward guidance with MCTS tree search for self-training, simultaneously
% improving policy and reward models.
% \autocite{guan2025rstarmathsmallllms} (rStar-Math) combine MCTS with
% a \emph{process preference model} (PPM) trained on step-level
% comparison data, improving Qwen2.5-Math-7B from 58.8\% to 90.0\% on
% MATH and matching o1-preview on competition benchmarks.  Both papers
% demonstrate that the self-evolving generator--verifier loop enabled by
% process rewards can close most of the gap with frontier closed models.

% % ============================================================
% \subsection{Reinforcement Learning with Verifiable Rewards}
% \label{sec:grpo}
% % ============================================================

% The most dramatic practical development since the original draft of
% this paper is the emergence of \emph{reasoning models} trained with
% RL using rule-based verifiable rewards, which has reshaped the
% context in which Critique DPO sits.

% \subsubsection{GRPO: Group Relative Policy Optimisation}

% \autocite{shao2024deepseekmath} (DeepSeekMath) introduce Group
% Relative Policy Optimisation (GRPO), which replaces PPO's learned
% critic with a group-score baseline, substantially reducing memory
% overhead:
% \begin{equation}
%   \hat{A}_{i,t}
%   = \frac{r_i - \mathrm{mean}(\{r_j\}_{j=1}^G)}
%          {\mathrm{std}(\{r_j\}_{j=1}^G)},
% \end{equation}
% where rewards are normalised within groups of $G$ sampled responses.
% GRPO has become the standard algorithm for reasoning model training.

% \subsubsection{DeepSeek-R1}

% \autocite{deepseekai2025deepseekr1} demonstrate that long-form
% CoT reasoning, self-reflection, and verification \emph{emerge} from
% pure RL with verifiable rewards, without any human-labelled reasoning
% trajectories.  DeepSeek-R1's four-stage pipeline
% (cold-start SFT~$\to$ RL~$\to$ rejection-sampling SFT~$\to$ final RL)
% and its distilled variants have become the de-facto baseline for
% open-source reasoning models.  The training paradigm is directly
% relevant to Critique DPO: R1 uses binary correctness as a reward,
% whereas Critique DPO proposes richer critique-structured rewards.

% \subsubsection{Kimi k1.5}

% \autocite{kimiteam2025kimik15} report that their \emph{Chain-of-Thought
% reward model} achieves 98.5\% annotation accuracy versus 84.4\% for a
% classical scalar reward model on a challenging preference dataset.
% This provides direct empirical support for the premise of Critique DPO:
% reward models that reason before judging are substantially more
% accurate than those that predict a scalar directly.

% \subsubsection{Open-Source GRPO Variants: DAPO and VAPO}

% \autocite{yu2025dapo} (DAPO) achieve 50 points on AIME 2024 with
% Qwen2.5-32B using asymmetric clipping, token-level normalisation,
% dynamic sampling, and an entropy bonus—surpassing DeepSeek-R1-Zero
% with 50\% fewer training steps.
% \autocite{liu2025vapo} (VAPO) additionally address value-model bias,
% heterogeneous sequence lengths, and reward sparsity in long-CoT RL,
% achieving 60.4 on AIME 2024.

% \subsubsection{Tulu 3 / RLVR}

% \autocite{lambert2024tulu3} demonstrate that an
% \emph{RLVR} (Reinforcement Learning with Verifiable Rewards) stage
% added after DPO produces the best open-source post-training recipe at
% the time of publication, providing an important practical template for
% pipelines that could incorporate Critique DPO signals.

% ============================================================
\subsection{Training Critique and Self-Correction Models}
\label{sec:critique-rl}
% ============================================================

A second growing cluster trains models specifically to \emph{generate}
critiques or corrections, rather than using externally provided ones.

\subsubsection{CriticGPT}

\autocite{mcaleese2024llmcriticshelpllm} train an LLM critique model
via RLHF on human-annotated errors injected into ChatGPT outputs.
CriticGPT-trained critics are preferred over unassisted human critics
in 63\% of cases and catch $\sim$85\% of bugs versus $\sim$25\% for
humans alone.  The paper validates the scalable-oversight motivation
for critique-based training and provides a practical data construction
recipe.

% \subsubsection{SCoRe: Self-Correction via RL}

% \autocite{kumar2024traininglanguagemodelsselfcorrect} (SCoRe)
% train LLMs to self-correct using a multi-turn online RL framework,
% addressing distribution mismatch and behaviour collapse that arise
% when supervised fine-tuning is used for self-correction.  SCoRe
% achieves 15.6\% gain on MATH and 9.1\% on HumanEval on Gemini models.
% The two-turn structure (generate~$\to$ critique~$\to$ revise) directly
% instantiates the critique loop that Critique DPO aims to train into
% the base model.

\subsubsection{CTRL: Teaching Critique via RL}

\autocite{xie2025teachinglanguagemodelscritique} (CTRL) train a
critic model with RL to maximise correction performance for a fixed
generator, without requiring human supervision.  The approach produces
critiques that serve as accurate generative reward signals and enables
test-time scaling through iterative critique–revision cycles, with up
to 106\% relative improvement on challenging code-generation
benchmarks.

\subsubsection{SuperCorrect}

\autocite{yang2024supercorrect} introduce a two-stage framework: (1)
hierarchical thought template distillation from a large teacher model,
and (2) cross-model collaborative DPO training on teacher-guided
error-correction traces.  SuperCorrect-7B surpasses DeepSeekMath-7B
by 7.8\%/5.3\% on MATH/GSM8K.  The error-correction trace is a
distilled form of the generative verifier output that Critique DPO
proposes to include in the DPO loss.

% \subsubsection{MultiCritique}

% \autocite{lan2024multicritic} address the limitation that single-LLM
% critiques—including those from GPT-4—contain systematic errors.
% MultiCritique aggregates feedback from multiple LLM agents during both
% SFT and RL stages to generate higher-quality critique training data,
% achieving GPT-4-level critique quality with a fine-tuned 7B model.

\subsubsection{Critique-GRPO}

\autocite{zhang2025critiquegrpoadvancingllm} is the closest
existing paper to Critique DPO's core proposal: it integrates
natural-language critiques \emph{and} numerical rewards jointly in a
GRPO-based RL framework.  The method enables LLMs to learn from both
initial responses and critique-guided self-refinements simultaneously,
achieving +4.4\% on Qwen2.5-7B-Base and +16.7\% pass@1 on AIME 2024.
The key difference from Critique DPO is the use of GRPO rather than
DPO, and the online RL rather than offline preference optimisation
setting.

\subsubsection{RL Tango: Joint Generator--Verifier Training}

\autocite{zha2025rltangoreinforcinggenerator} simultaneously train
both a policy generator and a generative process-level verifier using
interleaved RL, with the verifier rewarded solely on outcome
correctness without explicit process labels.  The co-evolving pair
achieves best-in-class performance on competition mathematics while
the verifier leads ProcessBench, providing practical evidence that the
joint generator--verifier training loop in Critique DPO can work
without pre-labelled critique data.

% ============================================================
\subsection{Generative Reward Models and Thinking LLM-as-Judge}
\label{sec:thinking-judge}
% ============================================================

A third cluster directly extends the \emph{generative verifier} idea
of \autocite{zhang2024generativeverifiersrewardmodeling} toward
preference ranking and reward modelling at scale.

\subsubsection{Prometheus 2 and RewardBench}

\autocite{kim2024prometheus2} release Prometheus 2, an open-source
evaluator LLM trained on a combined SFT + preference dataset that
supports both direct assessment and pairwise ranking using custom
scoring rubrics.  It achieves the highest correlation with human
evaluators among open-source alternatives, confirming the utility of
critique-style reasoning in evaluation.

\autocite{lambert2024rewardbench} introduce RewardBench, the standard
benchmark for evaluating reward models across chat, reasoning, and
safety domains.  The benchmark's \emph{reasoning} category is the most
relevant for evaluating Critique DPO-trained models, as it tests
preference models specifically on mathematical and coding tasks.

\subsubsection{EvalPlanner (Thinking LLM-as-a-Judge)}

\autocite{saha2025evalplanner} train an evaluator LLM to generate an
\emph{evaluation plan}, execute it through reasoning, and then deliver
a final judgement.  A self-training loop with synthetic preference
pairs iteratively refines the evaluation plans.  EvalPlanner achieves
93.9 on RewardBench—state-of-the-art among generative reward models
at the time of publication—demonstrating that planning-then-judging
substantially outperforms direct preference prediction.

\subsubsection{J1: Thinking LLM-as-a-Judge via RL}

\autocite{whitehouse2025j1incentivizingthinkingllmasajudge}
apply GRPO to teach judges to \emph{think before judging} by
converting judgement tasks into formats with verifiable rewards and
adding consistency-based reward bonuses to reduce positional bias.
J1-Qwen-32B outperforms o1-mini, o3, and DeepSeek-R1-671B on several
judge benchmarks, trained only on synthetic data.  The training
methodology is essentially Critique DPO's core idea applied to a
standalone judge model via RL.

\subsubsection{RM-R1: Reward Modelling as Reasoning}

\autocite{chen2025rmr1rewardmodelingreasoning}
train Reasoning Reward Models (ReasRMs) in two stages: (1)
chain-of-rubrics distillation from a large model, and (2) RL with
verifiable rewards.  For chat tasks the model generates evaluation
rubrics and justifications; for reasoning tasks it solves the problem
itself before judging.  RM-R1 outperforms INF-ORM-Llama3.1-70B and
GPT-4o on reward model benchmarks by up to 4.9\%, establishing that
structured reasoning before scalar prediction is the right paradigm
for reward modelling.

\subsubsection{DeepSeek-GRM and Self-Principled Critique Tuning (SPCT)}

\autocite{liu2025deepseekgrm} introduce SPCT, an online RL method
where the reward model adaptively generates \emph{principles and
  critiques} for each input before scoring, trained via RL to maximise
downstream alignment quality.  Inference-time scaling is achieved by
sampling multiple critique chains in parallel, guided by a
meta-reward model.  SPCT achieves better performance than
training-time scaling and avoids systematic biases present in
traditional reward models.  This can be viewed as a production-scale
implementation of the Critique DPO reward-modelling idea, with the
critic output used directly as the reward signal.

% \subsubsection{Think-J}

% \autocite{huang2025thinkjlearningthinkgenerative}
% build initial thinking ability in a judge model from curated reasoning
% traces, then optimise via RL using both offline critic-guided
% and online rule-based rewards.  The 32B model surpasses proprietary
% judges including Claude 3.5 Sonnet and GPT-4o without extra human
% annotations.  The critic-guided offline RL component mirrors the
% Critique DPO mechanism of using a critic's CoT to construct the
% preference signal.

% % ============================================================
% \subsection{DPO Theory and Iterative Variants}
% \label{sec:dpo-theory}
% % ============================================================

% \subsubsection{Reinforced Token Optimisation (RTO)}

% \autocite{zhong2024dpomeetspporeinforced} bridge DPO and PPO
% by extracting token-level reward functions from DPO's implicit reward
% and using them in PPO training.  RTO outperforms PPO by 7.5 points on
% AlpacaEval 2, validating the token-level credit assignment view of DPO
% that underlies step-level Critique DPO.

% \subsubsection{Bootstrapping with DPO Implicit Rewards (DICE)}

% \autocite{chen2024bootstrappinglanguagemodelsdpo} use DPO's implicit
% reward to generate preference datasets for subsequent DPO rounds
% without external feedback, achieving $>$8\% improvement in
% length-controlled win rate on AlpacaEval 2.  DICE demonstrates that
% the DPO implicit reward can self-improve through bootstrapping, an
% important property for the iterative training loop in Critique DPO.

% \subsubsection{Online Iterative RLHF}

% \autocite{dong2024rlhfworkflowrewardmodeling} provide a reproducible
% recipe for online iterative RLHF combining reward model training with
% online DPO, demonstrating that online feedback substantially
% outperforms offline DPO and providing a practical template for
% integrating critique-structured rewards.

% ============================================================
% NEW BibTeX ENTRIES
% Add these to refs.bib
% ============================================================

%% ---- Step-Level DPO ----

% @inproceedings{zhang2024chainpreferenceoptimization,
%   title   = {Chain of Preference Optimization: Improving
%               Chain-of-Thought Reasoning in {LLMs}},
%   author  = {Xuan Zhang and Chao Du and Tianyu Pang and
%               Qian Liu and Wei Gao and Min Lin},
%   booktitle = {NeurIPS},
%   year    = {2024},
%   eprint  = {2406.09136},
%   archivePrefix = {arXiv},
% }

% @misc{lu2024stepcontrolleddpoleveragingstepwise,
%   title   = {Step-Controlled {DPO}: Leveraging Stepwise Error
%               for Enhanced Mathematical Reasoning},
%   author  = {Zimu Lu and Aojun Zhou and Ke Wang and
%               Houxing Ren and Weikang Shi and Junting Pan and
%               Mingjie Zhan and Hongsheng Li},
%   year    = {2024},
%   eprint  = {2407.00782},
%   archivePrefix = {arXiv},
% }

% @misc{hwang2024selfexporeenhancingmathematical,
%   title   = {Self-Explore: Enhancing Mathematical Reasoning in
%               Language Models with Fine-grained Rewards},
%   author  = {Hyeonbin Hwang and Doyoung Kim and Seungone Kim and
%               Seonghyeon Ye and Minjoon Seo},
%   year    = {2024},
%   eprint  = {2404.10346},
%   archivePrefix = {arXiv},
% }

% @misc{yang2025thinkingpreferenceoptimization,
%   title   = {Thinking Preference Optimization},
%   author  = {Wang Yang and Hongye Jin and Jingfeng Yang and
%               Vipin Chaudhary and Xiaotian Han},
%   year    = {2025},
%   eprint  = {2502.13173},
%   archivePrefix = {arXiv},
% }

% @misc{gao2025uncoveringimpactchainofthought,
%   title   = {Uncovering the Impact of Chain-of-Thought Reasoning
%               for Direct Preference Optimization},
%   author  = {(multiple)},
%   year    = {2025},
%   eprint  = {2502.11656},
%   archivePrefix = {arXiv},
% }

% @misc{yuan2024eurus,
%   title   = {Advancing {LLM} Reasoning Generalists with
%               Preference Trees},
%   author  = {Lifan Yuan and Ganqu Cui and Hanbin Wang and
%               Ning Ding and others},
%   year    = {2024},
%   eprint  = {2404.02078},
%   archivePrefix = {arXiv},
% }

%% ---- Process Reward Models ----

% @inproceedings{lightman2023letsverifystepstep,
%   title   = {Let's Verify Step by Step},
%   author  = {Hunter Lightman and Vineet Kosaraju and Yura Burda
%               and others},
%   booktitle = {ICLR},
%   year    = {2024},
%   eprint  = {2305.20050},
%   archivePrefix = {arXiv},
% }

% @misc{setlur2024rewardingprogressscalingautomated,
%   title   = {Rewarding Progress: Scaling Automated Process
%               Verifiers for {LLM} Reasoning},
%   author  = {Amrith Setlur and Chirag Nagpal and Adam Fisch and
%               Xinyang Geng and Jacob Eisenstein and Rishabh Agarwal
%               and Alekh Agarwal and Jonathan Berant and
%               Aviral Kumar},
%   year    = {2024},
%   eprint  = {2410.08146},
%   archivePrefix = {arXiv},
% }

% @misc{yuan2024freeprocessrewardsprocess,
%   title   = {Free Process Rewards without Process Labels},
%   author  = {Lifan Yuan and Wendi Li and Huayu Chen and
%               Ganqu Cui and Ning Ding and Kaiyan Zhang and
%               Bowen Zhou and Zhiyuan Liu and Hao Peng},
%   year    = {2024},
%   eprint  = {2412.01981},
%   archivePrefix = {arXiv},
% }

% @misc{cui2025primereinforcementprocessimplicit,
%   title   = {Process Reinforcement through Implicit Rewards},
%   author  = {Ganqu Cui and Lifan Yuan and Zefan Wang and
%               Hanbin Wang and Yuchen Zhang and others},
%   year    = {2025},
%   eprint  = {2502.01456},
%   archivePrefix = {arXiv},
% }

% @misc{khalifa2025thinkprmprocessrewardmodels,
%   title   = {Process Reward Models That Think},
%   author  = {Muhammad Khalifa and Rishabh Agarwal and
%               Lajanugen Logeswaran and Jaekyeom Kim and Hao Peng
%               and Moontae Lee and Honglak Lee and Lu Wang},
%   year    = {2025},
%   eprint  = {2504.16828},
%   archivePrefix = {arXiv},
% }

% @misc{zhao2025genprm,
%   title   = {{GenPRM}: Scaling Test-Time Compute of Process
%               Reward Models via Generative Reasoning},
%   author  = {Jian Zhao and Runze Liu and Kaiyan Zhang and
%               Zhimu Zhou and others},
%   year    = {2025},
%   eprint  = {2504.00891},
%   archivePrefix = {arXiv},
% }

% @misc{zhang2025lessonsdevelopingprocessreward,
%   title   = {The Lessons of Developing Process Reward Models in
%               Mathematical Reasoning},
%   author  = {Zhenru Zhang and Chujie Zheng and Yangzhen Wu and
%               Beichen Zhang and Runji Lin and Bowen Yu and
%               Dayiheng Liu and Jingren Zhou and Junyang Lin},
%   year    = {2025},
%   eprint  = {2501.07301},
%   archivePrefix = {arXiv},
% }

% @inproceedings{zhang2024restmcts,
%   title   = {{ReST-MCTS*}: {LLM} Self-Training via Process Reward
%               Guided Tree Search},
%   author  = {Dan Zhang and Sining Zhoubian and Ziniu Hu and
%               Yisong Yue and Yuxiao Dong and Jie Tang},
%   booktitle = {NeurIPS},
%   year    = {2024},
%   eprint  = {2406.03816},
%   archivePrefix = {arXiv},
% }

% @misc{guan2025rstarmathsmallllms,
%   title   = {rStar-Math: Small {LLMs} Can Master Math Reasoning
%               with Self-Evolved Deep Thinking},
%   author  = {Xinyu Guan and Li Lyna Zhang and Yifei Liu and
%               Ning Shang and Youran Sun and Yi Zhu and Fan Yang
%               and Mao Yang},
%   year    = {2025},
%   eprint  = {2501.04519},
%   archivePrefix = {arXiv},
% }

%% ---- GRPO / Reasoning RL ----

% @misc{shao2024deepseekmath,
%   title   = {{DeepSeekMath}: Pushing the Limits of Mathematical
%               Reasoning in Open Language Models},
%   author  = {Zhihong Shao and Peiyi Wang and Qihao Zhu and
%               Runxin Xu and Junxiao Song and others},
%   year    = {2024},
%   eprint  = {2402.03300},
%   archivePrefix = {arXiv},
% }

% @misc{deepseekai2025deepseekr1,
%   title   = {{DeepSeek-R1}: Incentivizing Reasoning Capability in
%               {LLMs} via Reinforcement Learning},
%   author  = {{DeepSeek-AI}},
%   year    = {2025},
%   eprint  = {2501.12948},
%   archivePrefix = {arXiv},
% }

% @misc{kimiteam2025kimik15,
%   title   = {Kimi k1.5: Scaling Reinforcement Learning with
%               {LLMs}},
%   author  = {{Kimi Team}},
%   year    = {2025},
%   eprint  = {2501.12599},
%   archivePrefix = {arXiv},
% }

% @misc{yu2025dapo,
%   title   = {{DAPO}: An Open-Source {LLM} Reinforcement Learning
%               System at Scale},
%   author  = {Qiying Yu and Zheng Zhang and Ruofei Zhu and
%               Yufeng Yuan and others},
%   year    = {2025},
%   eprint  = {2503.14476},
%   archivePrefix = {arXiv},
% }

% @misc{liu2025vapo,
%   title   = {{VAPO}: Efficient and Reliable Reinforcement Learning
%               for Advanced Reasoning Tasks},
%   author  = {{Bytedance Seed Team}},
%   year    = {2025},
%   eprint  = {2504.05118},
%   archivePrefix = {arXiv},
% }

% @misc{lambert2024tulu3,
%   title   = {Tulu 3: Pushing Frontiers in Open Language Model
%               Post-Training},
%   author  = {Nathan Lambert and Jacob Morrison and
%               Valentina Pyatkin and others},
%   year    = {2024},
%   eprint  = {2411.15124},
%   archivePrefix = {arXiv},
% }

%% ---- Critique / Self-Correction ----

% @misc{mcaleese2024llmcriticshelpllm,
%   title   = {{LLM} Critics Help Catch {LLM} Bugs},
%   author  = {Nat McAleese and Rai Michael Pokorny and
%               Juan Felipe Ceron Uribe and Evgenia Nitishinskaya
%               and Maja Trebacz and Jan Leike},
%   year    = {2024},
%   eprint  = {2407.00215},
%   archivePrefix = {arXiv},
% }

% @misc{kumar2024traininglanguagemodelsselfcorrect,
%   title   = {Training Language Models to Self-Correct via
%               Reinforcement Learning},
%   author  = {Aviral Kumar and Vincent Zhuang and Rishabh Agarwal
%               and Yi Su and others},
%   year    = {2024},
%   eprint  = {2409.12917},
%   archivePrefix = {arXiv},
% }

% @misc{xie2025teachinglanguagemodelscritique,
%   title   = {Teaching Language Models to Critique via
%               Reinforcement Learning},
%   author  = {Zhihui Xie and Jie Chen and Liyu Chen and
%               Weichao Mao and Jingjing Xu and Lingpeng Kong},
%   year    = {2025},
%   eprint  = {2502.03492},
%   archivePrefix = {arXiv},
% }

% @misc{yang2024supercorrect,
%   title   = {{SuperCorrect}: Advancing Small {LLM} Reasoning with
%               Thought Template Distillation and Self-Correction},
%   author  = {Ling Yang and Zhaochen Yu and Tianjun Zhang and
%               Minkai Xu and Joseph E. Gonzalez and Bin Cui and
%               Shuicheng Yan},
%   year    = {2024},
%   eprint  = {2410.09008},
%   archivePrefix = {arXiv},
% }

% @misc{lan2024multicritic,
%   title   = {Training Language Models to Critique With
%               Multi-agent Feedback},
%   author  = {Tian Lan and Wenwei Zhang and Chengqi Lyu and
%               Shuaibin Li and Chen Xu and others},
%   year    = {2024},
%   eprint  = {2410.15287},
%   archivePrefix = {arXiv},
% }

% @misc{zhang2025critiquegrpoadvancingllm,
%   title   = {{Critique-GRPO}: Advancing {LLM} Reasoning with
%               Natural Language and Numerical Feedback},
%   author  = {Xiaoying Zhang and Hao Sun and Yipeng Zhang and
%               Kaituo Feng and Chaochao Lu and Chao Yang and
%               Helen Meng},
%   year    = {2025},
%   eprint  = {2506.03106},
%   archivePrefix = {arXiv},
% }

% @misc{zha2025rltangoreinforcinggenerator,
%   title   = {{RL Tango}: Reinforcing Generator and Verifier
%               Together for Language Reasoning},
%   author  = {Kaiwen Zha and Zhengqi Gao and Maohao Shen and
%               Zhang-Wei Hong and Duane S. Boning and Dina Katabi},
%   year    = {2025},
%   eprint  = {2505.15034},
%   archivePrefix = {arXiv},
% }

%% ---- Thinking Judges / Reasoning Reward Models ----

% @inproceedings{kim2024prometheus2,
%   title   = {Prometheus 2: An Open Source Language Model
%               Specialized in Evaluating Other Language Models},
%   author  = {Seungone Kim and Juyoung Suk and Shayne Longpre and
%               Bill Yuchen Lin and others},
%   booktitle = {EMNLP},
%   year    = {2024},
%   eprint  = {2405.01535},
%   archivePrefix = {arXiv},
% }

% @misc{lambert2024rewardbench,
%   title   = {{RewardBench}: Evaluating Reward Models for Language
%               Modeling},
%   author  = {Nathan Lambert and Valentina Pyatkin and
%               Jacob Morrison and LJ Miranda and Bill Yuchen Lin
%               and others},
%   year    = {2024},
%   eprint  = {2403.13787},
%   archivePrefix = {arXiv},
% }

% @misc{saha2025evalplanner,
%   title   = {Learning to Plan \& Reason for Evaluation with
%               Thinking-{LLM}-as-a-Judge},
%   author  = {Swarnadeep Saha and Xian Li and
%               Marjan Ghazvininejad and Jason Weston and
%               Tianlu Wang},
%   year    = {2025},
%   eprint  = {2501.18099},
%   archivePrefix = {arXiv},
% }

% @misc{whitehouse2025j1incentivizingthinkingllmasajudge,
%   title   = {J1: Incentivizing Thinking in {LLM}-as-a-Judge via
%               Reinforcement Learning},
%   author  = {Chenxi Whitehouse and Tianlu Wang and Ping Yu and
%               Xian Li and Jason Weston and Ilia Kulikov and
%               Swarnadeep Saha},
%   year    = {2025},
%   eprint  = {2505.10320},
%   archivePrefix = {arXiv},
% }

% @misc{chen2025rmr1rewardmodelingreasoning,
%   title   = {{RM-R1}: Reward Modeling as Reasoning},
%   author  = {Xiusi Chen and others},
%   year    = {2025},
%   eprint  = {2505.02387},
%   archivePrefix = {arXiv},
% }

% @misc{liu2025deepseekgrm,
%   title   = {Inference-Time Scaling for Generalist Reward
%               Modeling},
%   author  = {Zijun Liu and Peiyi Wang and Runxin Xu and
%               Shirong Ma and Chong Ruan and Peng Li and
%               Yang Liu and Yu Wu},
%   year    = {2025},
%   eprint  = {2504.02495},
%   archivePrefix = {arXiv},
% }

% @misc{huang2025thinkjlearningthinkgenerative,
%   title   = {Think-J: Learning to Think for Generative
%               {LLM}-as-a-Judge},
%   author  = {Hui Huang and Yancheng He and Hongli Zhou and
%               Rui Zhang and Wei Liu and Weixun Wang and
%               Jiaheng Liu and Wenbo Su},
%   year    = {2025},
%   eprint  = {2505.14268},
%   archivePrefix = {arXiv},
% }

%% ---- DPO Theory / Iterative ----

% @misc{zhong2024dpomeetspporeinforced,  % already in refs.bib
% }

% @misc{chen2024bootstrappinglanguagemodelsdpo,
%   title   = {Bootstrapping Language Models with {DPO} Implicit
%               Rewards},
%   author  = {Changyu Chen and Zichen Liu and Chao Du and
%               Tianyu Pang and Qian Liu and Arunesh Sinha and
%               Pradeep Varakantham and Min Lin},
%   year    = {2024},
%   eprint  = {2406.09760},
%   archivePrefix = {arXiv},
% }

% @misc{dong2024rlhfworkflowrewardmodeling,
%   title   = {{RLHF} Workflow: From Reward Modeling to Online
%               {RLHF}},
%   author  = {Hanze Dong and Wei Xiong and Bo Pang and
%               Haoxiang Wang and Han Zhao and Yingbo Zhou and
%               Nan Jiang and Doyen Sahoo and Caiming Xiong and
%               Tong Zhang},
%   year    = {2024},
%   eprint  = {2405.07863},
%   archivePrefix = {arXiv},
% }
